% main.tex -- RegRAG-VN, IEEE-RIVF 2026 (EDAS N35414)
%
% BUILD (verified on this machine, 2026-09-11):
%   cd paper && pdflatex main && bibtex main && pdflatex main && pdflatex main
% IEEEtran.cls V1.8b and IEEEtran.bst are NOT installed in this TeX Live tree;
% local copies (the genuine CTAN files, V1.8b, md5 5b2e4fa15b0f7eabb840ebf67df4c0f7)
% ship in this directory so the manuscript compiles standalone. Do not replace
% the document class with another one.
%
% NOTE(encoding): this host has no Vietnamese TeX support (no vntex, no babel
% `vietnam'), and lualatex is unusable (luaotfload broken). vnchars.sty supplies
% the missing precomposed characters as T1 composites and disables itself when
% a real Vietnamese setup is present. Compile on a full TeX Live for camera-ready.

\documentclass[conference,a4paper]{IEEEtran}
\ifdefined\pdfminorversion\pdfminorversion=6\fi

\usepackage[utf8]{inputenc}
\usepackage[T1]{fontenc}
\usepackage{vnchars}          % Vietnamese precomposed chars; see file header
\usepackage{cite}
\usepackage{amsmath,amssymb}
\usepackage{graphicx}
\usepackage{booktabs}
\usepackage{url}
\usepackage{xcolor}
\usepackage{tikz}
\usetikzlibrary{arrows.meta,positioning}

\begin{document}

\title{Citation-Level Unfaithfulness of Small Language Models on Vietnamese Banking Regulation: A Provenance-Guarded RAG Benchmark}

\author{%
\IEEEauthorblockN{%
Nguyen Huy Hoang\textsuperscript{1},
Dinh Thi Lan Huong\textsuperscript{2},
Nguyen Thang Phuc\textsuperscript{3},
Vu Duc Thanh\textsuperscript{4},
Nguyen Khac Duy Ngoc\textsuperscript{5}}
\IEEEauthorblockA{%
Hanoi University of Science and Technology, Hanoi, Vietnam\\
\textsuperscript{1}Hoang.NH20251325M@sis.hust.edu.vn \quad
\textsuperscript{2}Huong.DTL20261261M@sis.hust.edu.vn \quad
\textsuperscript{3}Phuc.NT20252263M@sis.hust.edu.vn\\
\textsuperscript{4}Thanh.VD20261082M@sis.hust.edu.vn \quad
\textsuperscript{5}Ngoc.NKD20261206M@sis.hust.edu.vn}
}

\maketitle

\begin{abstract}
Retrieval-augmented generation (RAG) is the standard remedy for hallucination in small language models, but its benefit is usually reported as an aggregate accuracy delta that hides where an answer came from. We present RegRAG-VN, a benchmark and evaluation framework for Vietnamese banking and payment regulation in which the retrieval corpus is aligned with the benchmark questions by requiring each question's supporting gold passage to match an ingested legal chunk exactly. The benchmark holds 64 answerable questions drawn from State Bank of Vietnam circulars, decrees, and laws, alongside 24 unanswerable probes paired with an explicit abstention sentinel. We evaluate models of at most 7B parameters under closed-book, BM25, and dense retrieval, scoring article-level (\emph{Điều}) and clause-level (\emph{Khoản}) citation accuracy alongside abstention. Across a $3\times3\times88$ generation campaign on a single GPU, retrieval augmentation raises article-level citation F1 from 0.005 (closed-book) to 0.32--0.43. However, 43--49\% of RAG-generated answers still cite an incorrect instrument or non-governing article, and models frequently emit the refusal sentinel before answering anyway. Finally, we describe a strict provenance discipline tracking corpus tiers and retriever backends, preventing degraded or fixture data from silently corrupting benchmark results.
\end{abstract}

\begin{IEEEkeywords}
retrieval-augmented generation, hallucination, citation faithfulness, abstention, Vietnamese legal NLP, benchmark
\end{IEEEkeywords}

\section{Introduction}
\label{sec:intro}

Regulatory compliance in Vietnam is a language problem with legal consequences. A bank officer who answers a customer's question about payment-card limits, or a compliance analyst who checks whether a new product violates a circular of the State Bank of Vietnam (SBV), must locate the governing provision among hundreds of circulars (\emph{Thông tư}), decrees (\emph{Nghị định}) and laws (\emph{Luật}), each organised into numbered articles (\emph{Điều}) and clauses (\emph{Khoản}). Language models are an obvious aid, but the models that regional banks and fintechs can actually deploy are small: a 3--7B parameter model served on a single commodity GPU, not a frontier API whose data leaves the country and whose pricing is set abroad. Whether such models answer Vietnamese regulatory questions faithfully and cite the governing provisions accurately is therefore a practical deployment question as well as a research one.

Retrieval-augmented generation is well established as a remedy for hallucination in English knowledge-intensive tasks \cite{lewis2020rag,gao2024ragsurvey}, and simply replicating that aggregate gain in Vietnamese provides limited practical insight. The failure mode that aggregate accuracy conceals is \emph{citation-level unfaithfulness}: a model produces a fluent, plausible answer that carries the wrong \emph{Điều} or \emph{Khoản}, or cites an instrument that does not govern the question at all, instead of abstaining. For a compliance tool, a wrong answer with a confident, wrong citation is more dangerous than an obviously wrong answer, because the citation is what a human reviewer checks first. Measuring that behaviour requires a benchmark in which the correct citation is known per question and in which the corpus provably contains the evidence, so that a missing citation cannot be blamed on missing evidence.

RegRAG-VN is designed around this requirement. Its corpus is \emph{QA-driven}: the legal instruments ingested are exactly those cited by the benchmark questions, and a coverage invariant requires every retained question's verbatim gold passage to string-match a chunk of the ingested text. Ground truth is therefore verifiably present in the retrieval corpus rather than assumed. On top of this corpus we measure three research questions:

\begin{itemize}
\item \textbf{RQ1 (hallucination reduction).} How much does RAG reduce hallucination versus closed-book answering for small models ($\leq$7B) on Vietnamese regulatory questions?
\item \textbf{RQ2 (citation and abstention).} How accurately do models cite the supporting article (\emph{Điều}) and clause (\emph{Khoản}), and do they abstain when the corpus does not contain the answer?
\item \textbf{RQ3 (sparse versus dense).} How does BM25 with Vietnamese compound-word segmentation compare with multilingual dense embeddings (BGE-M3) on legal retrieval?
\end{itemize}

A benchmark of this kind is only as trustworthy as the pipeline that produced its numbers. During development we found three code paths that silently substituted a degraded implementation and returned plausible-looking output; the most serious returned corpus-order chunks with fabricated similarity scores whenever the embedding library was absent, which would have invented an entire dense-retrieval column and invalidated RQ3 without any warning. We therefore treat provenance as a methodological requirement: every chunk, retrieved result, generation and evaluation record carries its corpus tier and retriever backend, and aggregation raises an exception rather than average any row whose provenance is unset, degraded, or a development fixture. Section~\ref{sec:method} describes this discipline, which ensures that silently fabricated columns cannot pass undetected as experimental results.

The contributions of this paper are: (i) RegRAG-VN, a QA-driven Vietnamese banking-regulation benchmark with 64 answerable questions, 24 unanswerable probes and a coverage invariant that makes evidence presence checkable; (ii) an evaluation protocol for citation-level faithfulness and abstention at \emph{Điều}/\emph{Khoản} granularity, including an evaluation of inter-annotator agreement on a cross-annotated subset; and (iii) a provenance discipline that makes silent substitution of degraded components impossible, documented against three concrete eliminated paths and one quarantined inauthentic seed corpus.

\section{Related Work}
\label{sec:related}

\textbf{RAG and its evaluation.} Retrieval-augmented generation conditions a generator on passages retrieved from a corpus \cite{lewis2020rag}, with dense passage retrieval \cite{karpukhin2020dpr} and sparse BM25 scoring \cite{robertson2009bm25} as the two dominant retriever families. Surveys organise the design space of retrievers, generators and evaluation \cite{gao2024ragsurvey}. Existing RAG benchmark suites evaluate component capabilities such as noise robustness, negative rejection, information integration, and counterfactual robustness \cite{liu2024rgb}, while automated frameworks estimate faithfulness and answer relevance without reference labels \cite{es2024ragas}. These suites are English-centric and none scores the correctness of a statutory citation, which is the unit a compliance reviewer trusts.

\textbf{Hallucination and faithfulness measurement.} Hallucination in natural language generation is surveyed along the intrinsic/extrinsic axis \cite{ji2023hallucination}. In summarisation, faithfulness was operationalised by evaluating whether source text entails summary facts \cite{maynez2020faithfulness}, while atomic-fact decomposition enabled precision scoring per claim \cite{min2023factscore}. Attribution research further formalises the requirement that generated statements be supported by cited sources, scoring citation quality separately from answer correctness \cite{rashkin2023attribution,bohnet2022attributed}. Self-critique architectures let models flag their own unsupported generations \cite{asai2024selfrag}. Our citation metric follows this attribution line but is stricter in one respect: the cited unit is a statutory article or clause, so a citation to the right instrument but the wrong article scores zero, matching how a legal reviewer reads a reference.

\textbf{Legal-domain language models.} Legal NLP benchmarks such as LexGLUE establish that legal text behaves differently from general text \cite{chalkidis2022lexglue}, and statute-specific retrieval has been studied for article-level matching. Vietnamese legal text adds a structural property that most legal-RAG work ignores: provisions are numbered hierarchically (\emph{Điều} within \emph{Chương}, \emph{Khoản} within \emph{Điều}), so the chunking granularity is itself an experimental variable. Prior work demonstrates that mixing chunk granularities outperforms any single fixed granularity for RAG \cite{qin2025mixgranularity}; our \emph{Điều}/\emph{Khoản} segmentation adapts this principle directly to the statutory hierarchy. Long-context dynamics are likewise relevant, as retrieved passages must compete for the model's attention within the prompt \cite{liu2024lostmiddle}.

\textbf{Vietnamese NLP.} Standard toolkits provide word segmentation and part-of-speech tagging for Vietnamese \cite{vu2018vncorenlp}. Segmentation is critical for sparse retrieval: because Vietnamese orthography places spaces between syllables rather than words, a whitespace tokenizer splits compound legal terms such as \emph{ngân hàng} or \emph{thẻ tín dụng} into fragmented units. We therefore apply word segmentation before BM25 scoring and evaluate the resulting segmented-sparse representation directly against multilingual dense embeddings (RQ3).

\section{The Benchmark}
\label{sec:benchmark}

\subsection{Domain and QA-driven corpus construction}
The domain is Vietnamese banking and payment regulation: SBV circulars on payment services, cards, non-cash payment and credit-institution organisation, together with the governing decrees and laws. Rather than querying a broad collection of legal instruments, the retrieval corpus is derived directly from the benchmark questions. Each gold question was authored together with a verbatim passage from the instrument that answers it and the source URL of that instrument. Instrument identifiers are resolved from the URL through a human-verified manifest; where a URL is opaque (a PDF slug or a document id that does not spell out the instrument number), the row is kept with an explicit \texttt{UNRESOLVED:} identifier and reported, never dropped, so that the benchmark cannot silently shrink. When a passage cites multiple candidate instruments, we designate the named instrument as primary, as spreadsheet URL ordering provides no substantive priority.

% NOTE(dataset): counts below verified 2026-09-11 against data/gold/bank_qa_data.csv
% via regrag/corpus/qa_loader.py; re-run scripts/load_qa_csv.py --report before
% camera-ready, the CSV is still under review.
The current benchmark holds 88 rows: 64 answerable questions, each with a gold passage recording the supporting provision and at least one source URL, and 24 unanswerable probes. Because gold passages reflect annotator records and may not represent verbatim extracts, the pipeline verifies each passage against the corpus at ingestion. When a recorded passage is a summary rather than an exact extract, the verbatim article text is substituted from the source instrument and logged; if the cited instrument is absent from the corpus, the row is excluded from the answerable set. In the delivered revision, 63 of the 64 passages consist of verbatim article text substituted under a guarded repair procedure. This procedure rewrites a row only if the cited instrument is present in the corpus, the cited article exists, and the annotator summary clears a lexical overlap threshold against that article (with all 63 substitutions logged alongside their originals). Of the remaining rows, one cites a \emph{Phụ lục} annex rather than a numbered article and is retained as an annotator summary, while another whose governing instrument is absent was transferred to the unanswerable probe arm. Since substituted passages originate directly from the ingested corpus, verbatim string coverage on repaired rows is partly self-satisfying; we therefore report it as an internal data-quality metric and anchor formal evaluation on article-level citation accuracy.
Article-level gold citations (\emph{Điều}) are derivable for 63 of the 64 answerable rows; clause-level gold (\emph{Khoản}) is present for 40. We therefore score citation accuracy at article level and treat clause level as optional: a clause reference is credited when present and correct but is never required for a match, because requiring it would penalise correct answers for questions whose governing provision has no clause structure.

\subsection{Two-tier corpus and the coverage invariant}
The corpus has two tiers with identical interfaces. Tier~1 is a development fixture built from the verbatim gold passages, one chunk per passage; it exists so that indexing, retrieval, prompting, generation and scoring could be integrated before the real instruments arrived. Because each question's gold passage is in Tier~1 by construction, retrieval recall on it is trivially perfect and sparse and dense retrievers are indistinguishable; Tier~1 is therefore never treated as corpus evidence: retrieval scores are not reported on it, and the generation campaign that used it does so only under the explicit disclosure of \S\ref{sec:disclosure}. Tier~2 is the segmented full text of the ingested instruments and is the corpus on which retrieval quality is measured (Table~\ref{tab:retrieval}). Segmentation follows the statutory hierarchy: chunks are articles (\emph{Điều}), with clause (\emph{Khoản}) boundaries preserved inside them, and preambles retained rather than discarded.

The coverage invariant enforces alignment between the tiers: every retained answerable question's verbatim gold passage must string-match a chunk of the Tier~2 corpus under strict normalisation (Unicode NFC plus whitespace collapse, diacritics preserved). A passage that matches only after diacritic folding is classified as extraction damage, as lost diacritics in the ingested text would corrupt both retrieval and citation scoring. Similarly, a passage matching only after removing all whitespace is logged as a second diagnostic category, identifying an artifact of official gazette PDFs that insert spurious intra-word spaces while retaining all characters and diacritics. Neither type of degradation is accepted as an exact match. Any question whose governing instrument could not be obtained is excluded from the answerable set; retaining it would falsely penalise models for hallucinating evidence that was never retrievable.

\subsection{Unanswerable probes}
The 24 probes evaluate the abstention component of RQ2 using questions deliberately unanswerable from the corpus; their gold target is the sentinel string \emph{Không có trong kho văn bản} (``not in the document corpus''). These scenarios represent plausible, in-domain banking situations whose governing provisions lie outside the repository. Examples include opening payment accounts for condominium management boards, kindergartens, and foreign trade-representative offices; collateral held by a spouse or an owner under 18; mortgage top-ups on property already securing another obligation; and the reporting frequency of the capital-adequacy ratio. The capital-adequacy question was originally annotated as answerable, but source verification revealed it is governed by an uningested statistical-reporting circular, prompting its reassignment to the probe arm.

By testing realistic questions governed by missing instruments, the probes evaluate abstention under a corpus-coverage gap rather than simple rejection of out-of-domain topics. This setting is intentionally demanding: models cannot rely on topical distance to trigger refusal, but must instead recognise that the retrieved evidence does not support the answer. We state as a limitation that the probe set contains no precision-trap and no cross-jurisdictional arm, so abstention on numerically near-miss questions and on foreign rules is not measured. A model that answers a probe fluently has failed abstention even if the answer is globally true, because the benchmark strictly evaluates what the corpus supports.

\subsection{Annotation and inter-annotator reliability}
Questions were authored by two domain annotators. Schedule constraints precluded dual annotation across all 88 rows, and computing Cohen's $\kappa$ \cite{cohen1960kappa} on self-graded items would measure only self-consistency. We therefore established a cross-annotated subset of 30 questions, fifteen authored by each annotator and graded independently by both; $\kappa$ is reported on this subset, with the remainder disclosed as single-annotator. This design is weaker than full dual annotation, and we state it as a limitation: half of the subset consists of each annotator's own questions, and single-annotator rows carry unmeasured author bias.

\section{Method}
\label{sec:method}

\subsection{Retrieval}
We compare three retrieval modes at top-$k=3$. Under \emph{closed-book}, the prompt contains only the question without context. For \emph{RAG-BM25}, queries and corpus chunks are segmented with PyVi \cite{pyvi} to preserve compound legal terms, followed by BM25 scoring \cite{robertson2009bm25,rankbm25}. For \emph{RAG-Dense}, queries and chunks are embedded using multilingual BGE-M3 \cite{chen2024m3} and ranked by cosine similarity. In both RAG configurations, the top-3 retrieved passages are concatenated into the prompt with identical formatting, isolating retrieval representation as the sole experimental variable.

\subsection{Generation}
We evaluate three instruction-tuned models of at most 7B parameters: Qwen2.5-7B-Instruct and Qwen2.5-3B-Instruct \cite{yang2024qwen25}, alongside the Vietnamese-focused Vistral-7B-Chat \cite{nguyen2023vistral} as an architecturally distinct third model. All models are served in 4-bit quantisation (NF4 weights, float16 compute) on a single 24\,GB GPU. Prior to evaluation, each model passes an automated memory preflight using worst-case prompt lengths, and prompt formatting is verified across sample outputs to confirm correct chat-template application. Every prompt includes explicit instructions specifying the sentinel refusal string, ensuring abstention is measured as a deliberate behavior rather than an omission.

\subsection{Evaluation metrics}
For each (question, model, mode) record we compute: (i) \emph{groundedness}, whether the answer's atomic claims are supported by the retrieved passages, following atomic-fact evaluation \cite{min2023factscore} and attribution scoring \cite{rashkin2023attribution}; (ii) \emph{hallucination}, a claim contradicted by the corpus or unsupported under closed-book; (iii) \emph{citation precision, recall and F1} at \emph{Điều} level, with \emph{Khoản} credited optionally, where a citation to the correct article of the wrong instrument scores zero; and (iv) \emph{abstention accuracy} on the probes, i.e.\ whether the model emits the sentinel and withholds a substantive answer. An earlier scorer version derived correctness entirely from citation overlap and flagged an uncited fabrication as non-hallucinated; that version is retained only as a tagged placeholder and is not used for any reported number. The replacement scorer is deterministic, lexical, and released in full; it is self-labelled unvalidated until human agreement is measured, and no number in \S\ref{sec:results} is human-adjusted. The release includes its validation protocol, specifying agreement on approximately 50 human-labelled responses plus Cohen's $\kappa$ on the cross-annotated subset of \S\ref{sec:benchmark}-D.

\subsection{Provenance discipline}
\label{sec:provenance}
Every chunk, retrieved result, generation and evaluation record carries two tags: \texttt{corpus\_source} (\texttt{tier1\_passages} or \texttt{tier2\_full}) and \texttt{retriever\_backend} (for example \texttt{bm25-rank\_bm25+pyvi} or \texttt{BAAI/bge-m3}, or \texttt{DEGRADED:<reason>}). The aggregation guard raises an exception rather than aggregating any row whose backend is unset or degraded or whose corpus is not Tier~2. Under this discipline, any placeholder must fail loudly or self-identify rather than return a plausible default.

Three silent-substitution paths motivated the rule. First, the dense index returned the first $k$ chunks in corpus order with fabricated scores whenever the embedding library was missing, which would have produced a fictional dense column and invalidated RQ3 without warning; it now raises a named error. Second, the sparse index silently degraded to term-overlap counting without its scoring library and to whitespace tokenisation without the Vietnamese segmenter, destroying exactly the compound-word handling RQ3 tests while still labelling itself BM25; the fallbacks remain but self-identify as degraded. Third, the evaluation scorers returned confident numbers from heuristics; they now tag themselves as placeholders.

A fourth incident concerned the corpus itself. A seed file that had become the entire corpus was audited and quarantined as inauthentic: 43 lines containing only articles 1, 2, 3, 9, 14, 15 and 23, whereas real instruments number articles contiguously; and lacking every mandatory element of a Vietnamese circular (recitals, effective-date clause, execution clause, signature block). Its chunks were removed and the segmenter, originally validated against that file, was re-validated against the real gold passages. Fabricated provisions are worse than no distractor, because a model can retrieve a fake article and be scored against it.

\section{Experimental Design and Results}
\label{sec:results}

We evaluate three models across three retrieval modes (closed-book, RAG-BM25, and RAG-Dense) at top-$k=3$ over all 88 benchmark rows, producing 792 generated answers. The campaign ran in a single window of about ten hours on a single RTX 3090 (24\,GB), serving each model in 4-bit NF4 quantisation with greedy decoding (temperature 0, seed 1234, at most 512 new tokens) and one model resident at a time. Every record carries its prompt hash, quantisation tag and GPU tag; all 792 prompts round-trip byte-identical through the harness, 709 answers finish naturally and 83 hit the token cap and are analysed as generated. Retrieval quality is reported separately as Recall@1 and Recall@3 on the answerable set (Table~\ref{tab:retrieval}), because a RAG arm whose retriever never finds the gold passage measures prompting rather than retrieval.

\subsection{Campaign corpus disclosure}
\label{sec:disclosure}
The generation campaign ran on the Tier-1 development fixture, not on Tier~2: the Tier-2 re-ingest was rejected by the harness's duplicate-chunk guard shortly before the deadline, and a post-re-ingest re-run did not fit the submission window. The fixture has one property that matters here: by construction, every answerable question's gold passage is in the corpus, so in both RAG arms the retrieved context provably contains the evidence. This setting carries three direct consequences. First, Tables~\ref{tab:main} and~\ref{tab:citation} measure the effect of \emph{conditioning on retrieved evidence}, not of retrieval quality; the RAG-vs-closed-book deltas are upper bounds. Second, on the fixture the BM25 and dense arms retrieve near-identical evidence sets, so their generation-level differences reflect prompt-packing rank order rather than retrieval quality; the two arms produced byte-identical answers for only 45 of 264 shared rows. Third, retrieval quality itself is measured separately on Tier~2 (Table~\ref{tab:retrieval}). Every row behind Tables~\ref{tab:main} and~\ref{tab:citation} is stamped \texttt{tier1\_passages} and is reported under this disclosure rather than as Tier-2 evidence, which the aggregation guard of \S\ref{sec:provenance} refuses.

% Table I: computed by scripts/compute_tables.py from data/eval/generations.json
% (log data/eval/tables_paper_2026-09-13.json); no cell hand-entered. Rows are
% tier1-passages-stamped and reported under the disclosure of \S\ref{sec:disclosure}.
\begin{table}[t]
\centering
\caption{Hallucination, groundedness and abstention by model and retrieval mode (RQ1, RQ2-abstention), from the 792-answer campaign of \S\ref{sec:disclosure}. Halluc.\ = hallucination rate over the 64 answerable questions; Ground.\ = mean sentence-level groundedness over rows with assessable retrieved context (closed-book has none, hence --); Abst.\ = share of the 24 probes correctly refused. All values are the deterministic scorer's output; no cell is hand-entered.}
\label{tab:main}
\begin{tabular}{@{}llccc@{}}
\toprule
Model & Mode & Halluc. $\downarrow$ & Ground. $\uparrow$ & Abst. $\uparrow$ \\
\midrule
Qwen2.5-7B  & closed-book & 0.438 & --    & 0.042 \\
Qwen2.5-7B  & RAG-BM25    & 0.344 & 0.985 & 0.500 \\
Qwen2.5-7B  & RAG-Dense   & 0.313 & 0.951 & 0.375 \\
Qwen2.5-3B  & closed-book & 0.750 & --    & 0.000 \\
Qwen2.5-3B  & RAG-BM25    & 0.469 & 0.950 & 0.000 \\
Qwen2.5-3B  & RAG-Dense   & 0.422 & 0.916 & 0.125 \\
Vistral-7B  & closed-book & 0.625 & --    & 0.000 \\
Vistral-7B  & RAG-BM25    & 0.672 & 0.912 & 0.000 \\
Vistral-7B  & RAG-Dense   & 0.594 & 0.939 & 0.000 \\
\bottomrule
\end{tabular}
\end{table}

% Table II: computed by scripts/compute_tables.py (log
% data/eval/tables_paper_2026-09-13.json); tier1-stamped rows under
% \S\ref{sec:disclosure}. 189 scorable answerable rows = 63 questions x 3
% models (Q036 cites a Phu luc annex, no article gold); clause denominator
% 120 = 40 questions x 3 models.
\begin{table}[t]
\centering
\caption{Citation accuracy at article (\emph{Điều}) level and optional clause (\emph{Khoản}) level, per retrieval mode, pooled over the three models (RQ2), from the campaign of \S\ref{sec:disclosure}. P and R are averaged per row over the 189 scorable answerable rows (63 questions $\times$ 3 models; one annex-citing question has no article gold); F1 is the harmonic mean of the pooled P and R; a correct article in the wrong instrument scores zero. \emph{Khoản} accuracy is over the 120 rows whose gold citation names a clause.}
\label{tab:citation}
\begin{tabular}{@{}lcccc@{}}
\toprule
Mode & \multicolumn{3}{c}{\emph{Điều} P / R / F1} & \emph{Khoản} acc. \\
\midrule
closed-book & 0.005 & 0.005 & 0.005 & 0.008 \\
RAG-BM25    & 0.293 & 0.354 & 0.321 & 0.300 \\
RAG-Dense   & 0.399 & 0.476 & 0.434 & 0.317 \\
\bottomrule
\end{tabular}
\end{table}

% Table III: BM25 row from scripts/eval_bm25_recall.py, log data/eval/
% bm25_recall_gate_c_2026-09-11.json (reproduced twice, 2026-09-11; reproduced
% a third time, identically, after the chunk-id fix). Dense row from the same
% script's dense arm, log data/eval/dense_recall_gate_c_2026-09-13.json
% (RunPod RTX 3090, 2026-09-13). doc-level = any retrieved chunk from a cited
% instrument, article-level = the cited Dieu among top-k. Article-level
% denominator is 63 (Q036 cites a Phu luc annex, no article id).
% TODO(results): kappa column from the 30-question cross-annotated subset
% (Thành & Hương, human).
\begin{table}[t]
\centering
\caption{Retrieval quality on the Tier~2 corpus and annotation agreement (RQ3, \S\ref{sec:benchmark}-D). Recall at instrument (doc) and article (\emph{Điều}) level, $k{=}1/3$, over the answerable set (doc $n{=}64$; article $n{=}63$, one row cites an annex); $\kappa$ is Cohen's kappa on the 30-question cross-annotated subset only.}
\label{tab:retrieval}
\setlength{\tabcolsep}{3pt}
\footnotesize
\begin{tabular}{@{}lccccc@{}}
\toprule
Retriever & Doc@1 & Doc@3 & Điều@1 & Điều@3 & $\kappa$ (subset) \\
\midrule
BM25 (pyvi) & 0.734 & 0.906 & 0.381 & 0.492 & -- \\
BGE-M3 dense & 0.719 & 0.891 & 0.365 & 0.524 & -- \\
\bottomrule
\end{tabular}
\end{table}

% Pipeline figure: TikZ, in-source, mirrors scripts/build_corpus_chunks.py ->
% scripts/run_campaign_pod.py -> scripts/compute_tables.py.
\begin{figure}[t]
\centering
\begin{tikzpicture}[
  node distance=1.1mm,
  box/.style={draw, rounded corners=1pt, font=\scriptsize, align=center,
              inner sep=2.2pt, minimum height=3.9mm,
              text width=0.88\columnwidth-6pt},
  arr/.style={-{Stealth[length=1.6mm]}}
]
\node[box] (qa)     {gold QA CSV $+$ source manifest \emph{(only trusted artifacts)}};
\node[box, below=of qa]     (ing) {QA-driven ingestion $\to$ Tier-2 corpus (\emph{Điều}/\emph{Khoản} chunks), Tier-1 fixture};
\node[box, below=of ing]    (cov) {coverage gate: verbatim gold passage $\wedge$ instrument present};
\node[box, below=of cov]    (idx) {BM25 (pyvi) / BGE-M3 indexing $\to$ top-3 context packing};
\node[box, below=of idx]    (gen) {3 models $\times$ 3 modes $\times$ 88 questions, 4-bit, greedy $\to$ 792 answers};
\node[box, below=of gen]    (agg) {lexical scoring + provenance-guarded aggregation \texttt{(corpus\_source, retriever\_backend)}};
\draw[arr] (qa)   -- (ing);
\draw[arr] (ing)  -- (cov);
\draw[arr] (cov)  -- (idx);
\draw[arr] (idx)  -- (gen);
\draw[arr] (gen)  -- (agg);
\end{tikzpicture}
\caption{The RegRAG-VN pipeline. Every record crossing a component boundary carries the \texttt{corpus\_source} and \texttt{retriever\_backend} tags enforced by the aggregation guard; the coverage gate excludes a question rather than let it generate a false hallucination label.}
\label{fig:pipeline}
\end{figure}

\subsection{RQ1: evidence conditioning reduces hallucination for two of three models}
Closed-book hallucination rates span 0.44--0.75 (Table~\ref{tab:main}). Conditioning on retrieved evidence lowers hallucination for both Qwen models (Qwen2.5-7B $0.438\!\to\!0.344/0.313$, Qwen2.5-3B $0.750\!\to\!0.469/0.422$ under BM25/dense). In contrast, Vistral-7B exhibits no reduction; its BM25 rate is marginally \emph{higher} than its closed-book baseline (0.672 vs.\ 0.625) because it paraphrases retrieved text into unsupported assertions rather than extracting it. Groundedness is high whenever context is present (0.91--0.98), as expected on a corpus where the evidence is verbatim in the prompt. The aggregate delta also masks an opposing calibration shift: Qwen2.5-7B falsely refuses 25.0\% and 23.4\% of answerable questions under RAG against 0\% closed-book, shifting the model from overconfident answering toward false refusal.

\subsection{RQ2: citation-level unfaithfulness is the dominant failure}
Closed-book citation performance is near zero: article-level F1 $= 0.005$, with 102 of 192 answered questions asserting a provision that does not govern the question and 76 citing an instrument other than the governing one (categories overlap). RAG lifts citation F1 to 0.321 (BM25) and 0.434 (dense), but mis-citation remains the modal failure: even with the governing provision inside the retrieved context, 49\% (BM25) and 43\% (dense) of answered rows still cite a wrong or non-covering provision, and clause-level (\emph{Khoản}) accuracy never exceeds 0.317. Model refusal is similarly brittle. Under closed-book evaluation, 71 of 72 probe rows receive substantive answers despite the sentinel instruction. Furthermore, 115 of 192 closed-book answers state the sentinel yet answer anyway; a simple keyphrase abstention scorer would credit these as correct refusals, motivating the assertion gate of \S\ref{sec:method}-C. Across all configurations, Vistral-7B refuses no probe in any mode.

Qualitative inspection of model answers reveals the underlying failure patterns. In question Q006 (closed-book), where the governing provision is Article~52(4) of the 2024 Notarization Law, Qwen2.5-7B instead cites ``\emph{Điều 405 Bộ luật Dân sự 2015}'' and appends a fabricated statutory quotation. In Q008 (closed-book), the model emits the required sentinel but immediately follows with a substantive explanation citing ``\emph{Nghị định 02/2021/NĐ-CP\ldots\ Điều 14, Điều 15}'', demonstrating the hedge-then-answer pattern on an irrelevant decree. In Q059 (RAG-BM25), even with the governing text of the 2024 Law on Credit Institutions present in the prompt, the model attributes Article~3 to Circular 61/2025/TT-NHNN---retaining the correct article number while substituting an unrelated circular. In regulatory compliance, such citation errors are particularly hazardous because human auditors use citations as their primary index for verification.

\subsection{RQ3: instrument-level retrieval is nearly solved, article-level is not}
On the Tier-2 corpus (Table~\ref{tab:retrieval}), BM25 with Vietnamese segmentation retrieves a chunk of the cited instrument for 90.6\% of answerable questions at $k{=}3$ (73.4\% at $k{=}1$), passing the $\geq 0.7$ retrieval gate, but places the cited \emph{Điều} itself in the top 3 for only 49.2\% (38.1\% at $k{=}1$). The dense arm passes the same gate (89.1\% doc-level at $k{=}3$) and reverses the article-level ranking at $k{=}3$ (52.4\% against BM25's 49.2\%), while BM25 maintains a slight instrument-level lead (90.6\% vs.\ 89.1\%). The two retrievers thus exhibit opposing relative strengths across the two granularities. The substantial gap between instrument-level and article-level retrieval persists across both retrievers (a factor of 1.8$\times$ for BM25 and 1.7$\times$ for dense). While retrieving the governing circular or law is largely solved at $k{=}3$, placing the specific governing article in context remains difficult; this retrieval bottleneck bounds the citation ceiling of downstream generation, consistent with the mis-citation rates observed in RQ2.

% TODO(results): two or three further narrative points may be added if page
% budget allows: latency spread (Vistral RAG rows 164--175 s vs.\ Qwen 5--9 s)
% and the 83 token-capped answers.

\section{Conclusion}
\label{sec:conclusion}

RegRAG-VN exposes failure modes that aggregate accuracy metrics conceal by verifying evidence presence in the corpus, scoring citations at statutory article and clause resolution, and measuring model abstention when governing provisions are absent. Conditioning small models on retrieved evidence reduces unsupported answers for both Qwen models but not for Vistral-7B. It lifts article-level citation F1 from a closed-book floor of 0.005 to 0.32--0.43, yet 43--49\% of answered questions still carry confident citations to a wrong instrument or non-covering provision despite evidence in context. Clause-level accuracy remains at or below 0.32 across all configurations. Refusal behaviour proves equally brittle: Vistral-7B never abstains on any probe, while Qwen2.5-7B falsely refuses up to a quarter of answerable questions once evidence is present. On the full Tier-2 corpus, document-level retrieval at $k{=}3$ is largely solved for both retrievers (BM25 Recall@3 $= 0.906$, BGE-M3 $0.891$), whereas article-level recall falls to $0.492$ (BM25) and $0.524$ (BGE-M3). Because the two retrievers swap relative rank across granularities, this bottleneck reflects statutory hierarchy rather than the sparse-versus-dense choice. The provenance discipline proved essential in practice: development uncovered and eliminated three silent-substitution paths and one inauthentic seed corpus, and the duplicate-chunk guard that rejected the Tier-2 re-ingest enforced the Tier-1 disclosure of \S\ref{sec:disclosure} rather than allowing a fixture to pass as corpus evidence. Limitations qualify these findings: the generation campaign ran on the Tier-1 fixture under disclosure; lexical scorers remain deterministic and self-labelled unvalidated pending completion of the human-agreement protocol (approximately 50 labelled responses and the 30-question cross-annotated $\kappa$ subset), with 58 of 88 rows single-annotator; and evaluation spans three $\leq$7B models under 4-bit quantisation on one GPU. Corpus, code, scorers, prompts and the full 792-row campaign log are released for reproduction.

\section*{Data and Code Availability}
% TODO(results): insert the Zenodo DOI here at submission; the record must be
% published before the EDAS entry is finalised.
The benchmark, ingestion and evaluation code, provenance guard and all campaign logs will be released as a public repository with a Zenodo DOI at submission; the placeholder for the DOI is reserved in this paragraph. The trusted gold CSV, instrument manifest and coverage reports serve as the artifacts of record; every derived file is regenerable from them by a single idempotent script. Legal texts are redistributed as downloaded from official government sources and remain under their original terms.

\bibliographystyle{IEEEtran}
\bibliography{refs}

\end{document}
